\clearpage

\section{Порядок работы системы}

Запрос: \texttt{?- find\_grandparent(GP, female, female, ivan).} (Поиск бабушки по материнской линии).

\begin{tabularx}{\textwidth}{|c|>{\small}X|>{\small}X|>{\small}X|}
	\hline
	\textbf{№}                                                                         &
	\textbf{Состояние резольвенты}                                                     &
	\textbf{Для каких термов запускается алгоритм унификации, результат и подстановка} &
	\textbf{Дальнейшие действия}                                                                 \\
	\hline

	0                                                                                  &
	find\_grandparent(GP, female, female, ivan).                                       &
	—                                                                                  &
	Резольвента непустая. \newline
	Прямой ход. \newline
	Запуск алгоритма унификации для верхней подцели с 1-ым заголовком из БЗ (find\_grandparent). \\
	\hline

	1                                                                                  &
	parent(GP, Parent), \newline
	gender(GP, GPGender), \newline
	parent(Parent, Child), \newline
	gender(Parent, PGender). \newline
	Применение подстановки: \newline
	parent(GP, Parent), \newline
	gender(GP, female), \newline
	parent(Parent, ivan), \newline
	gender(Parent, female).                                                            &
	find\_grandparent(GP, female, female, ivan) =
	find\_grandparent(GP, GPGender, PGender, Child) \newline
	Унификация успешна. \newline
	Подстановка: \newline
	\{GPGender = female, \newline PGender = female, \newline Child = ivan\}            &
	Унификация успешна. \newline
	Резольвента непустая. \newline
	Прямой ход.                                                                                  \\
	\hline

	2                                                                                  &
	parent(GP, Parent), \newline
	gender(GP, female), \newline
	parent(Parent, ivan), \newline
	gender(Parent, female).                                                            &
	parent(GP, Parent) = \newline
	find\_grandparent(...) \newline
	Унификация неуспешна (разные функторы).                                            &
	Унификация неуспешна. \newline
	Резольвента непустая. \newline
	Прямой ход.                                                                                  \\
	\hline

	3                                                                                  &
	parent(GP, Parent), \newline
	gender(GP, female), \newline
	...                                                                                &
	parent(GP, Parent) = \newline
	gender(petr, male) \newline
	Унификация неуспешна (разные функторы).                                            &
	Унификация неуспешна. \newline
	Резольвента непустая. \newline
	Прямой ход.                                                                                  \\
	\hline

	4-9                                                                                &
	...                                                                                &
	...                                                                                &
	... (Пропуск оставшихся 6 фактов gender)                                                     \\
	\hline

	10                                                                                 &
	parent(GP, Parent), \newline
	gender(GP, female), \newline
	parent(Parent, ivan), \newline
	gender(Parent, female). \newline
	Редукция и применение подстановки к полученной конъюнкции целей: \newline
	gender(petr, female), \newline
	parent(maria, ivan), \newline
	gender(maria, female).                                                             &
	parent(GP, Parent) = \newline
	parent(petr, maria) \newline
	Унификация успешна. \newline
	Подстановка: \newline
	\{GP = petr, Parent = maria\}                                                      &
	Унификация успешна. \newline
	Выполнение редукции подцели и применение подстановки.
	Резольвента непустая. \newline
	Прямой ход.                                                                                  \\
	\hline

	11                                                                                 &
	gender(petr, female), \newline
	parent(maria, ivan), \newline
	gender(maria, female).                                                             &
	gender(petr, female) = \newline
	find\_grandparent(...) \newline
	Унификация неуспешна (разные функторы).                                            &
	Унификация неуспешна. \newline
	Резольвента непустая. \newline
	Прямой ход.                                                                                  \\
	\hline

	12                                                                                 &
	gender(petr, female), \newline
	parent(maria, ivan), \newline
	gender(maria, female).                                                             &
	gender(petr, female) = \newline
	gender(petr, male) \newline
	Унификация неуспешна (разные константы).                                           &
	Унификация неуспешна. \newline
	Резольвента непустая. \newline
	Прямой ход.                                                                                  \\
	\hline

	13-20                                                                              &
	...                                                                                &
	...                                                                                &
	... (Пропуск остальных фактов БЗ. Везде неуспех)                                             \\
	\hline

	21                                                                                 &
	gender(petr, female), \newline
	parent(maria, ivan), \newline
	gender(maria, female). \newline
	Восстановленное состояние резольвенты: \newline
	parent(GP, Parent), \newline
	gender(GP, female), \newline
	parent(Parent, ivan), \newline
	gender(Parent, female).                                                            &
	Встречен конец БЗ.                                                                 &
	Откат. Восстанавливается предыдущее состояние резольвенты (Шаг 10). \newline
	Резольвента непустая. \newline
	Прямой ход. \newline
	Запуск алгоритма унификации для parent(GP, Parent) и 10-ого заголовка parent(anna, maria).   \\
	\hline

	22                                                                                 &
	parent(GP, Parent), \newline
	gender(GP, female), \newline
	parent(Parent, ivan), \newline
	gender(Parent, female). \newline
	Успешная унификация. \newline
	Редукция и применение подстановки к полученной конъюнкции целей: \newline
	gender(anna, female), \newline
	parent(maria, ivan), \newline
	gender(maria, female).                                                             &
	parent(GP, Parent) = \newline
	parent(anna, maria) \newline
	Унификация успешна. \newline
	Подстановка: \newline
	\{GP = anna, Parent = maria\}                                                      &
	Унификация успешна. \newline
	Резольвента непустая. \newline
	Прямой ход.                                                                                  \\
	\hline

	23-26                                                                              &
	...                                                                                &
	...                                                                                &
	... (Пропуск правила find\_grandparent и первых фактов gender)                               \\
	\hline

	27                                                                                 &
	gender(anna, female), \newline
	parent(maria, ivan), \newline
	gender(maria, female). \newline
	Успешная унификация. \newline
	Редукция и применение подстановки к полученной конъюнкции целей: \newline
	parent(maria, ivan), \newline
	gender(maria, female).                                                             &
	gender(anna, female) = \newline
	gender(anna, female) \newline
	Унификация успешна. \newline
	Подстановка: \newline
	Нет.                                                                               &
	Унификация успешна. \newline
	Резольвента непустая. \newline
	Прямой ход.                                                                                  \\
	\hline

	28-39                                                                              &
	...                                                                                &
	...                                                                                &
	... (Пропуск несовпадающих заголовков с 1 по 12)                                             \\
	\hline

	40                                                                                 &
	parent(maria, ivan), \newline
	gender(maria, female). \newline
	Успешная унификация. \newline
	Редукция и применение подстановки к полученной конъюнкции целей: \newline
	gender(maria, female).                                                             &
	parent(maria, ivan) = \newline
	parent(maria, ivan) \newline
	Унификация успешна. \newline
	Подстановка: \newline
	Нет.                                                                               &
	Унификация успешна. \newline
	Резольвента непустая. \newline
	Прямой ход.                                                                                  \\
	\hline

	41-47                                                                              &
	...                                                                                &
	...                                                                                &
	... (Пропуск несовпадающих заголовков с 1 по 7)                                              \\
	\hline

	48                                                                                 &
	gender(maria, female). \newline
	Успешная унификация. \newline
	Редукция и применение подстановки к полученной конъюнкции целей: \newline
	Нет.                                                                               &
	gender(maria, female) = \newline
	gender(maria, female) \newline
	Унификация успешна. \newline
	Подстановка: \newline
	Нет.                                                                               &
	Унификация успешна. \newline
	Резольвента пуста. \newline
	Знания подобраны. Подстановка \{GP = anna\}. \newline
	Откат (при запросе следующего решения).                                                      \\
	\hline
\end{tabularx}

% \clearpage
Запрос: \texttt{?- max3\_no\_cut(0, 5, 2, Res).}

\begin{tabularx}{\textwidth}{|c|>{\small}X|>{\small}X|>{\small}X|}
	\hline
	\textbf{№}                                                                         &
	\textbf{Состояние резольвенты}                                                     &
	\textbf{Для каких термов запускается алгоритм унификации, результат и подстановка} &
	\textbf{Дальнейшие действия}                                                                 \\
	\hline

	0                                                                                  &
	max3\_no\_cut(0, 5, 2, Res).                                                       &
	—                                                                                  &
	Резольвента непустая. \newline
	Прямой ход. \newline
	Запуск алгоритма унификации для верхней подцели с 1-ым заголовком из БЗ (find\_grandparent). \\
	\hline

	1                                                                                  &
	max3\_no\_cut(0, 5, 2, Res).                                                       &
	max3\_no\_cut(0, 5, 2, Res) = \newline
	find\_grandparent(...) \newline
	Унификация неуспешна (разные функторы).                                            &
	Унификация неуспешна. \newline
	Резольвента непустая. \newline
	Прямой ход.                                                                                  \\
	\hline

	2-23                                                                               &
	...                                                                                &
	... (Пропуск несовпадающих правил gender, parent, max2\_no\_cut, max2\_cut)        &
	...                                                                                          \\
	\hline

	24                                                                                 &
	max3\_no\_cut(0, 5, 2, Res). \newline
	Успешная унификация. \newline
	Редукция и применение подстановки к полученной конъюнкции целей: \newline
	>=(0, 5), \newline
	>=(0, 2).                                                                          &
	max3\_no\_cut(0, 5, 2, Res) = max3\_no\_cut(X, Y, Z, X) \newline
	Унификация успешна. \newline
	Подстановка: \newline
	\{X = 0, Y = 5, Z = 2, Res = X\}                                                   &
	Унификация успешна. \newline
	Резольвента непустая. \newline
	Прямой ход. \newline
	Запуск алгоритма унификации для >=(0, 5) и Истина.                                           \\
	\hline

	25                                                                                 &
	>=(0, 5), \newline
	>=(0, 2). \newline
	Восстановленное состояние резольвенты: \newline
	max3\_no\_cut(0, 5, 2, Res).                                                       &
	>=(0, 5) = \newline
	Ложь \newline
	Унификация неуспешна (Ложь).                                                       &
	Унификация неуспешна. \newline
	Откат. Восстанавливается предыдущее состояние резольвенты (Шаг 0). \newline
	Резольвента непустая. \newline
	Прямой ход.                                                                                  \\
	\hline

	26                                                                                 &
	max3\_no\_cut(0, 5, 2, Res). \newline
	Успешная унификация. \newline
	Редукция и применение подстановки к полученной конъюнкции целей: \newline
	>(5, 0), \newline
	>=(5, 2).                                                                          &
	max3\_no\_cut(0, 5, 2, Res) = max3\_no\_cut(X, Y, Z, Y) \newline
	Унификация успешна. \newline
	Подстановка: \newline
	\{X = 0, Y = 5, Z = 2, Res = Y\}                                                   &
	Унификация успешна. \newline
	Резольвента непустая. \newline
	Прямой ход. \newline
	Запуск алгоритма унификации для >(5, 0) и Истина.                                            \\
	\hline

	27                                                                                 &
	>(5, 0), \newline
	>=(5, 2). \newline
	Успешная унификация. \newline
	Редукция и применение подстановки к полученной конъюнкции целей: \newline
	>=(5, 2).                                                                          &
	>(5, 0) = \newline
	Истина \newline
	Унификация успешна. \newline
	Подстановка: \newline
	Нет.                                                                               &
	Унификация успешна. \newline
	Резольвента непустая. \newline
	Прямой ход. \newline
	Запуск алгоритма унификации для >=(5, 2) и Истина.                                           \\
	\hline

	28                                                                                 &
	>=(5, 2). \newline
	Успешная унификация. \newline
	Редукция и применение подстановки к полученной конъюнкции целей: \newline
	Нет.                                                                               &
	>=(5, 2) = \newline
	Истина \newline
	Унификация успешна. \newline
	Подстановка: \newline
	Нет.                                                                               &
	Унификация успешна. \newline
	Резольвента пуста. \newline
	Знания подобраны. Подстановка \{Res = 5\}. \newline
	Откат (при запросе следующего решения).                                                      \\
	\hline

	29                                                                                 &
	max3\_no\_cut(0, 5, 2, Res). \newline
	Успешная унификация. \newline
	Редукция и применение подстановки к полученной конъюнкции целей: \newline
	>(2, 0), \newline
	>(2, 5).                                                                           &
	max3\_no\_cut(0, 5, 2, Res) = max3\_no\_cut(X, Y, Z, Z) \newline
	Унификация успешна. \newline
	Подстановка: \newline
	\{X = 0, Y = 5, Z = 2, Res = Z\}                                                   &
	Унификация успешна. \newline
	Резольвента непустая. \newline
	Прямой ход. \newline
	Запуск алгоритма унификации для >(2, 0) и Истина.                                            \\
	\hline

	30                                                                                 &
	>(2, 0), \newline
	>(2, 5). \newline
	Успешная унификация. \newline
	Редукция и применение подстановки к полученной конъюнкции целей: \newline
	>(2, 5).                                                                           &
	>(2, 0) = \newline
	Истина \newline
	Унификация успешна. \newline
	Подстановка: \newline
	Нет.                                                                               &
	Унификация успешна. \newline
	Резольвента непустая. \newline
	Прямой ход. \newline
	Запуск алгоритма унификации для >(2, 5) и Истина.                                            \\
	\hline

	31                                                                                 &
	>(2, 5). \newline
	Восстановленное состояние резольвенты: \newline
	max3\_no\_cut(0, 5, 2, Res).                                                       &
	>(2, 5) = \newline
	Ложь \newline
	Унификация неуспешна (Ложь).                                                       &
	Унификация неуспешна. \newline
	Откат. Восстанавливается предыдущее состояние резольвенты (Шаг 0). \newline
	Резольвента непустая. \newline
	Прямой ход.                                                                                  \\
	\hline

	32-34                                                                              &
	...                                                                                &
	... (Пропуск оставшихся правил max3\_cut)                                          &
	...                                                                                          \\
	\hline

	35                                                                                 &
	max3\_no\_cut(0, 5, 2, Res).                                                       &
	Встречен конец БЗ.                                                                 &
	Откат. К предыдущему состоянию. \newline
	Резольвента пуста. \newline
	Завершение работы.                                                                           \\
	\hline
\end{tabularx}

\clearpage
Запрос: \texttt{?- max3\_cut(0, 5, 2, Res).}

\begin{tabularx}{\textwidth}{|c|>{\small}X|>{\small}X|>{\small}X|}
	\hline
	\textbf{№}                                                                                     &
	\textbf{Состояние резольвенты}                                                                 &
	\textbf{Для каких термов запускается алгоритм унификации, результат и подстановка}             &
	\textbf{Дальнейшие действия}                                                                     \\
	\hline

	0                                                                                              &
	max3\_cut(0, 5, 2, Res).                                                                       &
	—                                                                                              &
	Резольвента непустая. \newline
	Прямой ход. \newline
	Запуск алгоритма унификации для верхней подцели с 1-ым заголовком из БЗ (find\_grandparent).     \\
	\hline

	1                                                                                              &
	max3\_cut(0, 5, 2, Res).                                                                       &
	max3\_cut(0, 5, 2, Res) = \newline
	find\_grandparent(...) \newline
	Унификация неуспешна (разные функторы).                                                        &
	Унификация неуспешна. \newline
	Резольвента непустая. \newline
	Прямой ход.                                                                                      \\
	\hline

	2-23                                                                                           &
	...                                                                                            &
	... (Пропуск несовпадающих предикатов gender, parent, max2\_no\_cut, max2\_cut, max3\_no\_cut) &
	...                                                                                              \\
	\hline

	24                                                                                             &
	max3\_cut(0, 5, 2, Res). \newline
	Успешная унификация. \newline
	Редукция и применение подстановки к полученной конъюнкции целей: \newline
	>=(0, 5), \newline
	>=(0, 2), \newline
	!.                                                                                             &
	max3\_cut(0, 5, 2, Res) = \newline
	max3\_cut(X, Y, Z, X) \newline
	Унификация успешна. \newline
	Подстановка: \newline
	\{X = 0, Y = 5, Z = 2, Res = X\}                                                               &
	Унификация успешна. \newline
	Резольвента непустая. \newline
	Прямой ход. \newline
	Запуск алгоритма унификации для >=(0, 5) и Истина.                                               \\
	\hline

	25                                                                                             &
	>=(0, 5), \newline
	>=(0, 2), \newline
	!. \newline
	Восстановленное состояние резольвенты: \newline
	max3\_cut(0, 5, 2, Res).                                                                       &
	>=(0, 5) = \newline
	Ложь \newline
	Унификация неуспешна (Ложь).                                                                   &
	Унификация неуспешна. \newline
	Откат. Восстанавливается предыдущее состояние резольвенты (Шаг 0). \newline
	Резольвента непустая. \newline
	Прямой ход.                                                                                      \\
	\hline

	26                                                                                             &
	max3\_cut(0, 5, 2, Res). \newline
	Успешная унификация. \newline
	Редукция и применение подстановки к полученной конъюнкции целей: \newline
	>=(5, 2), \newline
	!.                                                                                             &
	max3\_cut(0, 5, 2, Res) = \newline
	max3\_cut(\_, Y, Z, Y) \newline
	Унификация успешна. \newline
	Подстановка: \newline
	\{Y = 5, Z = 2, Res = Y\}                                                                      &
	Унификация успешна. \newline
	Резольвента непустая. \newline
	Прямой ход. \newline
	Запуск алгоритма унификации для >=(5, 2) и Истина.                                               \\
	\hline

	27                                                                                             &
	>=(5, 2), \newline
	!. \newline
	Успешная унификация. \newline
	Редукция и применение подстановки к полученной конъюнкции целей: \newline
	!.                                                                                             &
	>=(5, 2) = \newline
	Истина \newline
	Унификация успешна. \newline
	Подстановка: \newline
	Нет.                                                                                           &
	Унификация успешна. \newline
	Резольвента непустая. \newline
	Прямой ход. \newline
	Запуск алгоритма обработки предиката отсечения (!).                                              \\
	\hline

	28                                                                                             &
	!. \newline
	Успешная унификация. \newline
	Редукция и применение подстановки к полученной конъюнкции целей: \newline
	Нет.                                                                                           &
	! = \newline
	! \newline
	Унификация успешна. \newline
	Подстановка: \newline
	Нет.                                                                                           &
	Унификация успешна. \newline
	Предикат отсечения срабатывает: удаляются все точки возврата для текущей цели. \newline
	Резольвента пуста. \newline
	Знания подобраны. Подстановка \{Res = 5\}. \newline
	Завершение работы (альтернативные ветви отсечены).                                               \\
	\hline
\end{tabularx}
